% !TEX program = xelatex
% 本科毕业论文（PMFS-KV）——严格按常见高校毕设版式：A4、1.5 倍行距、中英文摘要、目录、正文章节、参考文献、附录、致谢
\documentclass[UTF8,a4paper,zihao=-4,oneside,openany]{ctexbook}

% ---------- 页面与行距 ----------
\usepackage[a4paper,left=2.5cm,right=2.5cm,top=2.5cm,bottom=2.5cm,
  headheight=15pt]{geometry}
\usepackage{setspace}
\onehalfspacing
\usepackage{fancyhdr}
\usepackage{titlesec}
\usepackage{titletoc}
\usepackage{booktabs}
\usepackage{tabularx}
\usepackage{longtable}
\usepackage{multirow}
\usepackage{makecell}
\usepackage{graphicx}
\usepackage{float}
\usepackage{caption}
\usepackage{subcaption}
\usepackage{hyperref}
\usepackage{xcolor}
\usepackage{listings}
\usepackage{amsmath,amssymb}
\usepackage{enumitem}
\usepackage{siunitx}
\usepackage{cleveref}

\hypersetup{
  colorlinks=true,
  linkcolor=black,
  citecolor=black,
  urlcolor=blue!50!black,
  bookmarksnumbered=true,
  pdftitle={面向持久内存的文件系统设计与实现——基于xv6的PMFS文件子系统},
  pdfauthor={本科毕业论文}
}

\captionsetup{font=small,labelsep=quad}
\captionsetup[figure]{position=bottom}
\captionsetup[table]{position=top}

% 字体：Noto CJK（WSL 已装）
\setCJKmainfont{Noto Serif CJK SC}[
  BoldFont=Noto Serif CJK SC Bold,
  AutoFakeSlant=0.2
]
\setCJKsansfont{Noto Sans CJK SC}[
  BoldFont=Noto Sans CJK SC Bold
]
\setCJKmonofont{Noto Sans Mono CJK SC}
\setmainfont{TeX Gyre Termes}
\setsansfont{TeX Gyre Heros}
\setmonofont{TeX Gyre Cursor}

% 章节标题
\titleformat{\chapter}{\centering\bfseries\Large}{第\,\thechapter\,章}{1em}{}
\titleformat{\section}{\bfseries\large}{\thesection}{1em}{}
\titleformat{\subsection}{\bfseries\normalsize}{\thesubsection}{1em}{}
\titlespacing*{\chapter}{0pt}{12pt}{18pt}

% 页眉页脚
\pagestyle{fancy}
\fancyhf{}
\fancyhead[C]{\small 面向持久内存的文件系统设计与实现——xv6 PMFS v3}
\fancyfoot[C]{\thepage}
\renewcommand{\headrulewidth}{0.4pt}

% 代码列表
\lstset{
  basicstyle=\ttfamily\footnotesize,
  breaklines=true,
  frame=single,
  columns=fullflexible,
  keepspaces=true,
  showstringspaces=false,
  keywordstyle=\bfseries,
  commentstyle=\itshape\color{gray!60!black},
}

% 精修实验图（pgfplots）——数据来自 docs/BENCH.md / make bench-sweep（CDC 融合后）
\usepackage{pgfplots}
\pgfplotsset{compat=1.18,
  every axis/.append style={
    font=\small,
    label style={font=\small},
    tick label style={font=\footnotesize},
    legend style={font=\footnotesize},
    grid=both,
    grid style={line width=.1pt, draw=gray!25},
    major grid style={line width=.2pt,draw=gray!40},
  }
}

% 图：对数轴 N–延迟（量级对比）
\newcommand{\FigNLatency}{%
\begin{figure}[htbp]
  \centering
  \begin{tikzpicture}
  \begin{axis}[
    width=0.92\textwidth,
    height=0.52\textwidth,
    xlabel={$N$（每 trial 键/文件数）},
    ylabel={中位数延迟（$r\_time$ 相对单位，对数轴）},
    ymode=log,
    xmin=4, xmax=36,
    ymin=1e4, ymax=5e7,
    xtick={8,16,32},
    legend pos=north west,
    legend cell align=left,
  ]
  \addplot+[thick, mark=*, mark size=2.5pt] coordinates {
    (8,33507) (16,38448) (32,53065)
  };
  \addlegendentry{PMFS-KV}
  \addplot+[thick, mark=triangle*, mark size=2.2pt, blue!50!black] coordinates {
    (8,304451) (16,486211) (32,659056)
  };
  \addlegendentry{PMFS-File}
  \addplot+[thick, dashed, mark=square*, mark size=2.2pt, red!70!black] coordinates {
    (8,5684939) (16,11208338) (32,25648605)
  };
  \addlegendentry{xv6 fs+sync}
  \end{axis}
  \end{tikzpicture}
  \caption{随 $N$ 变化的中位数延迟（EMUL，5 trial，CDC 融合后）。对数轴强调相对 xv6 的量级差，不作同等 POSIX 加速比解读。}
  \label{fig:n-latency}
\end{figure}
}

% 图：WA vs N
\newcommand{\FigWaN}{%
\begin{figure}[htbp]
  \centering
  \begin{tikzpicture}
  \begin{axis}[
    width=0.85\textwidth,
    height=0.48\textwidth,
    xlabel={$N$},
    ylabel={$\mathrm{wa\_x100}=\mathrm{phys}\times100/\mathrm{logical}$},
    xmin=4, xmax=36,
    ymin=0, ymax=22000,
    xtick={8,16,32},
  ]
  \addplot+[thick, mark=*, mark size=2.5pt, blue!60!black] coordinates {
    (8,15384) (16,17889) (32,17721)
  };
  \node at (axis cs:20,6000) [align=left, font=\footnotesize] {
    小 value（8\,B）相对 4\,KiB 页\\
    对齐主导放大；融合后随 $N$\\
    近似平台，属页粒度预期
  };
  \end{axis}
  \end{tikzpicture}
  \caption{写放大 $\mathrm{wa\_x100}$ 随 $N$ 变化（EMUL，KV 路径）。高放大属页粒度预期，与 xv6 日志放大不同源。}
  \label{fig:wa-n}
\end{figure}
}

% 图：MAP vs EMUL 柱状（KV 路径）
\newcommand{\FigMapEmul}{%
\begin{figure}[htbp]
  \centering
  \begin{tikzpicture}
  \begin{axis}[
    ybar,
    width=0.75\textwidth,
    height=0.45\textwidth,
    ylabel={KV 中位数延迟（$r\_time$）},
    symbolic x coords={EMUL,MAP},
    xtick=data,
    ymin=0, ymax=140000,
    bar width=18pt,
    nodes near coords,
    every node near coord/.append style={font=\footnotesize},
  ]
  \addplot coordinates {(EMUL,53065) (MAP,113804)};
  \end{axis}
  \end{tikzpicture}
  \caption{$N{=}32$ 时 PMFS-KV 在 EMUL 与 MAP 下的中位数延迟（CDC 融合后）。}
  \label{fig:map-emul}
\end{figure}
}

% 图：文件路径相对吞吐
\newcommand{\FigThr}{%
\begin{figure}[htbp]
  \centering
  \begin{tikzpicture}
  \begin{axis}[
    ybar,
    width=0.85\textwidth,
    height=0.48\textwidth,
    ylabel={相对吞吐（$N\times10^{6}/\mathrm{median}$）},
    symbolic x coords={N8,N16,N32},
    xtick=data,
    ymin=0, ymax=60,
    bar width=12pt,
    legend style={at={(0.98,0.98)},anchor=north east},
  ]
  \addplot coordinates {(N8,26) (N16,32) (N32,48)};
  \addplot coordinates {(N8,1) (N16,1) (N32,1)};
  \legend{PMFS-File,xv6 fs+sync}
  \end{axis}
  \end{tikzpicture}
  \caption{EMUL 下文件 API 相对吞吐（合题「吞吐」指标；非墙钟 MB/s）。}
  \label{fig:thr}
\end{figure}
}


% ---------- 封面信息（可按学校模板替换）----------
\newcommand{\ThesisTitle}{面向持久内存的文件系统设计与实现}
\newcommand{\ThesisSubtitle}{基于 xv6 的 PMFS v3：受限\texttt{/pm}、COW 提交与可复现验收}
\newcommand{\ThesisAuthor}{（姓名）}
\newcommand{\ThesisAdvisor}{（指导教师）}
\newcommand{\ThesisSchool}{（学院）}
\newcommand{\ThesisMajor}{计算机科学与技术}
\newcommand{\ThesisDate}{二〇二六\,年\,六\,月}

\begin{document}

% ========== 封面 ==========
\begin{titlepage}
\thispagestyle{empty}
\centering
\vspace*{1.2cm}
{\Large\bfseries 本科毕业论文\par}
\vspace{1.8cm}
{\zihao{2}\bfseries \ThesisTitle\par}
\vspace{0.6cm}
{\large \ThesisSubtitle\par}
\vspace{2.2cm}
\begin{tabular}{rl}
专业名称： & \ThesisMajor \\[0.6em]
学生姓名： & \ThesisAuthor \\[0.6em]
指导教师： & \ThesisAdvisor \\[0.6em]
所在学院： & \ThesisSchool \\[0.6em]
完成日期： & \ThesisDate \\
\end{tabular}
\vfill
{\large （按学校正式封面模板替换本页）\par}
\end{titlepage}

\frontmatter
\pagenumbering{Roman}

% ========== 中文摘要 ==========
\chapter*{摘\quad 要}
\addcontentsline{toc}{chapter}{摘要}
\markboth{摘要}{摘要}

持久内存（Persistent Memory, PMEM）以字节寻址提供近似 DRAM 的访问延迟并在掉电后保留数据，对以块设备与日志为中心的操作系统存储栈提出了新的一致性模型与验收问题。教学操作系统 xv6 以 virtio 与日志型文件系统为根，适合讲授经典 UNIX 语义，但难以直接承载「声明假设下的可复现提交验收」。

本文按原题四要件，在 xv6-riscv 上实现 \textbf{PMFS v3}：保留根文件系统，经固定挂载点 \texttt{/pm} 提供受限文件子系统；以 COW 元数据与不可变 KV run 组织字节寻址布局，使文件数据路径不经 \texttt{bio}/\texttt{log}/\texttt{virtio}；以行级 flush 请求、可选 Zicbom \texttt{CBO.CLEAN} 与双槽 checkpoint 提交词实现轻量提交协议；为 QEMU 9.2.0 提供实验性独立 NVDIMM 补丁路径。语义上写死可见点与 POSIX 差集：仅成功 \texttt{fsync}/checkpoint 后承诺完整版本可读回。

核心贡献是教学 OS 上\textbf{可复现的提交验收方法学}（进程重启切点 oracle、有界规格负控、轨迹采样）及受限但可测的文件子系统实现。实验分两套口径：架构面报告对齐工作量下相对原生 \texttt{fs+log} 的延迟/短批吞吐；公平面报告同契约软件请求消融；并明确设备模型与持久语义不等价。

\textbf{关键词：}持久内存；NVDIMM；文件系统；崩溃一致性；写放大；xv6；契约测试；COW

% ========== 英文摘要 ==========
\chapter*{Abstract}
\addcontentsline{toc}{chapter}{Abstract}
\markboth{Abstract}{Abstract}

Persistent memory (PMEM) offers byte-addressable access with near-DRAM latency while retaining data across power loss, challenging OS storage stacks built around block devices and journaling. The teaching OS xv6 provides a virtio-backed log file system as root; it is excellent for classic UNIX semantics but does not by itself provide reproducible acceptance of PM commit protocols under explicit assumptions.

This thesis implements \textbf{PMFS v3}, a bounded persistent-memory file subsystem on a fixed \texttt{/pm} mount beside the xv6 root: COW metadata and immutable KV runs without the block I/O stack; line-granular flush requests with optional Zicbom \texttt{CBO.CLEAN}; dual-slot checkpoints; and an experimental QEMU 9.2.0 NVDIMM patch path. Visibility is stated explicitly: a successful \texttt{fsync} publishes a complete version; unsynchronized updates are not promised after failure.

The distinctive contribution is a \textbf{reproducible commit-acceptance methodology} (process-kill cut-point oracles, bounded specification negatives, sampled traces) plus a measurable teaching-OS file subset. Evaluation separates an architectural plane (matched-workload latency/throughput versus xv6 \texttt{fs+log}, with unequal persist contracts) from a fairness plane (same-PMFS software-request ablations).

\textbf{Keywords:} persistent memory; NVDIMM; file system; crash consistency; write amplification; xv6; contract testing; copy-on-write

\tableofcontents
\clearpage

\mainmatter
\pagenumbering{arabic}

\chapter{绪论}
\label{chap:intro}

\section{研究背景}

操作系统存储子系统长期以块设备为中心：文件内容经缓冲缓存、日志与块 I/O 栈落盘，崩溃一致性依赖有序刷写与检查点。教学操作系统 xv6~\cite{xv6book} 以 virtio-blk 与日志型文件系统讲授这一经典路径，语义清晰、代码可阅读，但延迟受块设备与元数据更新约束，且难以自然表达字节寻址持久内存（Persistent Memory, PMEM）上的提交与掉电语义。

PMEM 以字节寻址提供近似 DRAM 的访问延迟，并在掉电后保留数据~\cite{nvmm2021}。学术界与工业界已出现面向 PMEM 的文件系统与编程模型，如 PMFS~\cite{pmfs2014}、NOVA~\cite{nova2016}、BPFS~\cite{bpfs2009}、Mnemosyne~\cite{mnemosyne2011} 等。它们多面向通用或接近生产的 Linux 环境；将完整 POSIX VFS 移植进 xv6 既超出本科毕设合理体量，也易掩盖「如何验收崩溃正确性」这一更尖锐的工程问题。合题要求的「文件子系统 + 刷写/NVDIMM 仿真 + 吞吐」，需要的是\textbf{可陈述、可测的子集语义}，而非声称完整生产 FS。

\section{问题陈述}

本文关注三个相互耦合的问题：

\begin{enumerate}[leftmargin=2em]
  \item \textbf{形态选择：}在 xv6 上如何提供并行于 virtio 根的 PMEM \textbf{文件子系统}（目录/extent + \texttt{create}/\texttt{fsync}），同时以 KV/run 核承载切分后的内容块，而不替换根 FS、也不假装完整 VFS？
  \item \textbf{一致性模型：}双槽超级块、不可变 run、批量 seal、inode 页 magic+校验，如何共同保证「未 \texttt{fsync}/未提交不可见、半写活动槽与半页元数据可检测」？
  \item \textbf{验收方法：}如何让崩溃、remount 字节等价与故障注入语义\textbf{可机器判定}（串口一行、宿主脚本、CI），而非仅靠人工读日志？
\end{enumerate}

\section{研究内容与目标}

本文设计并实现 \textbf{PMFS}：并行于 virtio 根的 PMEM 文件子系统（\texttt{create}/\texttt{mkdir}/\texttt{fopen}/\texttt{fread}/\texttt{fwrite}/\texttt{fsync}/\texttt{readdir}/\texttt{funlink}）与融合 KV 引擎（\texttt{put}/\texttt{get}/\texttt{seal}）；写入经 CDC/stream 切分后落入同一 run 核；后端支持 EMUL、MAP 与 PMEM2。研究目标：

\begin{itemize}[leftmargin=2em]
  \item 映射 QEMU NVDIMM 仿真窗；inode/dentry 去块设备布局；行级刷写协议；
  \item 建立含 \texttt{PERSIST}/\texttt{FILE}/\texttt{REMOUNT}/\texttt{FAULT}/\texttt{PMEM2} 的 \texttt{ONBOX\_PMFS\_*} 契约验收门；
  \item 报告文件 API 延迟中位数与相对吞吐，并以 A0--A4 公平消融给出切分/刷写效应量；
  \item 写死崩溃可见点与 POSIX 差集（见语义清单），使答辩口径可复核。
\end{itemize}

\section{主要贡献}

对照相关工作与作者前序原生项目（FCS / Aurora），本文贡献收敛为一条尖刺与两条支撑柱~\cite{fcsown,auroraown,cxcos}：

\begin{description}[leftmargin=2em,style=nextline]
  \item[尖刺（C1）] 教学 OS 上可机器判定的崩溃与掉电契约套件（含 remount 字节等价与故障注入）：相对同类，不可替代的是方法学验收，而非又一个双槽。
  \item[系统柱（C2）] 去块设备文件布局 + CDC/stream 融合 KV 核 + 行级 persist + 双槽/B3/WA。
  \item[实验柱（C3）] PMEM2 掉电一行过 + A0--A4 公平消融效应量（\texttt{cycles}/\texttt{wa\_x100}）。
\end{description}

主动不宣称的内容包括：完整内核 VFS、真 NVDIMM 主线化、将 FCS/Aurora 源码并入库、多写者并发、\texttt{lseek}/\texttt{mmap}/\texttt{rename}。

\section{论文结构}

第~\ref{chap:related}~章综述相关工作与借鉴边界；第~\ref{chap:design}~章给出总体设计（含可见点与差集）；第~\ref{chap:impl}~章描述实现与契约；第~\ref{chap:eval}~章报告实验、消融、性能图与踩坑；第~\ref{chap:concl}~章总结并展望。附录给出真 NVDIMM 与复现命令。

\chapter{相关工作}
\label{chap:related}

\section{教学操作系统中的块文件系统}

xv6~\cite{xv6book} 的 \texttt{fs+log} 路径以 inode、目录项、位图与日志检查点保证崩溃后可恢复。该设计适合教学，但其 I/O 路径绑定 virtio 块设备与 bio，元数据更新成本随「每文件一次 create/write/sync」近似线性增长。本文将该路径作为对照基线，\textbf{仍保留为根文件系统}，不主张用 PMFS-KV 替换日常用户程序的文件语义。

\section{持久内存文件系统与编程模型}

PMFS~\cite{pmfs2014} 与 NOVA~\cite{nova2016} 面向 PMEM 提供文件语义；Strata~\cite{strata2017} 等探索跨介质分层。BPFS~\cite{bpfs2009}、Mnemosyne~\cite{mnemosyne2011} 强调细粒度持久与编程抽象。上述工作的问题域是通用或接近生产的 PMEM 文件/运行时。本文取其「字节寻址 + 显式持久」动机，但产品形态收敛为 \textbf{KV 子集}，体量与 xv6 教学代码库匹配。

\section{日志结构与近似索引}

LSM 类结构~\cite{blsm2012} 与 Bloom 过滤器~\cite{bloom1970} 广泛用于写入友好与存在性过滤。本文 run 不可变、批量 seal 与 Bloom/fence 属于同一设计家族，但实现服务于教学内核约束（页粒度、有限段数、无用户态复杂压缩管线）。

\section{作者前序原生工作与工程契约手法}

作者前序原生项目 \textbf{FCS}~\cite{fcsown} 提供嵌入式耐久 KV 的设计约束：双槽超级块、不可变 run、B1/B2/B3 与写放大计量；\textbf{Aurora}~\cite{auroraown} 提供键量化、二值 posting 与查询 LRU 等索引约束。二者代码不并入本仓库，xv6 侧为 C 语言自研重写，以防「移植题」误读。

工程验收上，作者在其他内核工程中使用过 \texttt{ONBOX\_*} 串口契约手法~\cite{cxcos}；本课题将其重写为服务存储语义的 \texttt{ONBOX\_PMFS\_*} 阶段/失败码，使崩溃与掉电可被 \texttt{grep} 与 CI 判定。

\section{本文定位小结}

\begin{table}[htbp]
  \centering
  \caption{相关工作与本文关系}
  \label{tab:related}
  \begin{tabular}{p{3.2cm}p{4.2cm}p{5.2cm}}
    \toprule
    工作 & 问题域 & 与本文 \\
    \midrule
    xv6 fs+log & 教学块 FS & 对照基线；仍作根 FS \\
    Linux PMFS/NOVA 等 & PMEM 文件 FS & 动机参考；本文为 KV 子集 \\
    作者 FCS / Aurora & 耐久引擎与索引（原生前序） & 约束下沉；代码自研重写 \\
    cxc-os 类契约 & 串口阶段验收 & 手法参考；码表重写 \\
    \textbf{本文 PMFS-KV} & 教学 OS 上 PMEM-KV + 契约套件 & \textbf{尖刺 = 可机器判定崩溃/掉电} \\
    \bottomrule
  \end{tabular}
\end{table}

防「移植题」的标准表述：数据结构思想可追溯前序约束，但验收资产（契约套件）与 xv6 内实现为本课题自有贡献。

\chapter{系统设计}
\label{chap:design}

\section{产品形态与非目标}

PMFS v3 的产品形态是：\textbf{保留 virtio 根，经固定合成挂载点 \texttt{/pm} 提供受限文件子系统}。用户可用普通 \texttt{open}/\texttt{read}/\texttt{write}/\texttt{lseek}/\texttt{fsync}/\texttt{getdents} 等访问 PM 文件；并保留专用 KV/stream 接口。\textbf{不}替换根 FS，\textbf{不}做动态 mount table 或完整 POSIX。

非目标：上游通用 NVDIMM/NFIT/热插拔；无平台假设的物理掉电证明；符号链接、权限、\texttt{mmap}、跨后端 rename、打开后删除最后链接；全设备 GC/硬件磨损均衡。

\section{地址空间与后端}

\begin{lstlisting}[caption={地址布局示意（MAP / 实验 NVDIMM）},label={lst:layout}]
KERNBASE ---------- 0x80000000
  kernel + user（易失 DRAM，典型 128MiB）
PHYSTOP  ---------- KERNBASE + 128MiB
PMEMBASE ---------- PHYSTOP（如 0x88000000）
  PMFS 窗口：双槽 checkpoint / 元数据 / run
PMEMEND  ---------- PMEMBASE + 32MiB（名义；MAP 受 DTB 截断）
\end{lstlisting}

\begin{table}[htbp]
  \centering
  \caption{物理后端与可验证范围}
  \label{tab:backends}
  \begin{tabular}{p{2.2cm}p{5.0cm}p{5.2cm}}
    \toprule
    模式 & 含义 & 不可据此声称 \\
    \midrule
    EMUL & 8\,MiB 静态易失缓冲 & 跨启动保留、硬件持久 \\
    普通 MAP & \texttt{PHYSTOP} 后 guest RAM & 独立 NVDIMM 设备 \\
    共享 RAM 文件 & file-backend，便于杀进程重启 & 宿主掉电、独立设备 \\
    实验 NVDIMM & QEMU 9.2.0 补丁冷插拔 \texttt{nvdimm}+私有 FDT
      & 上游通用 virt、NFIT、物理掉电 \\
    Zicbom（可选） & 发出 \texttt{CBO.CLEAN} & TCG 下的真实缓存写回 \\
    \bottomrule
  \end{tabular}
\end{table}

MAP 在窗口不足或设备缺失时\textbf{拒绝挂载}，不静默降级为 EMUL。

\section{软件分层与对象图}

自上而下：VFS 逐组件遍历（原生 inode 与 PM inode pin）$\rightarrow$ \texttt{pmfile}（共享工作缓冲、COW 元数据）$\rightarrow$ \texttt{pmstream}/\texttt{pmcdc} $\rightarrow$ \texttt{pmalgo}/\texttt{pmindex} $\rightarrow$ \texttt{pmfs}（双槽 checkpoint、可达性、页复用）$\rightarrow$ \texttt{pmdev}/\texttt{pmem}。

文件内容以 KV chunk（extent）承载；元数据为不可变快照页。容量边界（量级）：约 40 inode/dentry、每文件 4 extent、文件 $\le$1024\,B。

\section{崩溃可见点与 POSIX 差集}
\label{sec:semantics}

\begin{table}[htbp]
  \centering
  \caption{提交可见性（软件协议）}
  \label{tab:visibility}
  \begin{tabular}{p{3.4cm}p{2.4cm}p{6.5cm}}
    \toprule
    操作 & 故障后承诺 & 说明 \\
    \midrule
    未成功 \texttt{fsync} 的写 & \textbf{不承诺}保留 & 工作缓冲/memtable \\
    成功 \texttt{fsync} & 完整新版本可读回 & 数据+元数据+提交词 \\
    空白/v2/校验损坏 & 拒绝挂载 & 不自动 format \\
    单槽损坏 & 可回退另一槽 & 双槽均坏则失败 \\
    \bottomrule
  \end{tabular}
\end{table}

差集包括：动态挂载、完整跨目录/跨后端 rename、符号链接、权限、\texttt{mmap}、多写者线性化证明等。硬链接、\texttt{lseek}/\texttt{O\_APPEND}、PM 内子树 rename、\texttt{getdents} 已支持（见实现章）。

\section{刷写与 checkpoint 发布顺序}

持久化请求语义依赖平台：默认 \texttt{pmem\_flush} 为按行 touch+\texttt{fence}（可观测、可排序）；可选 Zicbom 发出 \texttt{CBO.CLEAN}。一次根发布的设计顺序为：

\begin{enumerate}[leftmargin=2em]
  \item 追加 run / 分配 COW 元数据页；对各页请求 persist + fence；
  \item 非当前槽：清提交词并 persist；写入新头（未提交）并 persist；
  \item \textbf{一次 8\,B store}写入提交状态与校验 $\rightarrow$ persist + fence；再更新内存 epoch/槽。
\end{enumerate}

该协议还依赖：屏障顺序、提交词故障原子性、映射有效、校验未碰撞等假设。QEMU TCG 不模拟缓存持久域，故正文将「掉电安全」表述为\textbf{声明假设下的协议正确性}，并以进程重启 oracle 验证。

\section{维护、CDC 与评测意图}

双槽可达性保护两代根；双槽均不可达页可复用。显式 \texttt{pmcompact} 做全链 COW 压缩，近满或缺暂存空间可失败。普通写优先 CDC，超预算则固定分块；\textbf{无内容去重}。

评测分两平面：\textbf{架构面}——相对 xv6 的同工作量延迟/吞吐（持久契约不等价）；\textbf{公平面}——同 PMFS 内软件请求/策略消融。禁止混读。

\chapter{系统实现}
\label{chap:impl}

\section{模块组织}

实现基于 xv6-riscv。主要模块包括：VFS 路径遍历与 \texttt{/pm} 挂载点；\texttt{pmfile}（普通 fd、COW 元数据、硬链接/rename）；\texttt{pmcdc}/\texttt{pmstream}；\texttt{pmalgo}/\texttt{pmindex}；\texttt{pmfs}（checkpoint、可达性、压缩）；\texttt{pmdev}/\texttt{pmem}；以及 \texttt{tools/pmfs\_recovery.py}、\texttt{test-xv6.py}、契约与 bench 脚本。文件内容经 \texttt{pm\_put}/\texttt{seal} 落入 run；元数据经 COW 页与 checkpoint 发布。

\section{VFS 与普通文件路径}

\texttt{vfs\_walk} 逐组件解析；仅原生根下的 \texttt{pm} 进入 PMFS，\texttt{/pm/..} 返回原生根。PM inode pin 支持 cwd/目录 fd；独立 \texttt{open} 各自维护 offset，共享工作缓冲；\texttt{dup}/\texttt{fork} 共享打开对象。PMFS sleep lock 串行化缓冲、append 取尾与同步。

硬链接维护共享 inode 与链接数；非最后链接可在打开时删除，最后链接的 unlink-open 仍拒绝。子树 rename 检查环与路径长度后，修改内存图并以\textbf{一次} checkpoint 发布；非替换可复用不可变 inode 页。

\section{COW 元数据、run 与 seal}

内容同步：CDC/固定切分 $\rightarrow$ \texttt{pm\_put} $\rightarrow$ seal 追加不可变 run $\rightarrow$ 新 inode 快照（可复用未变 dentry）$\rightarrow$ checkpoint。\texttt{pm\_get} 优先 memtable/缓存，再精确索引或 newest-first + Bloom/fence。索引满后未跟踪键走 fallback；更新已跟踪键必须刷新指向。

安全挂载：无有效/旧格式/损坏则失败关闭；仅 \texttt{pmformat --force} 显式格式化。v2 与 v3 不兼容，无迁移工具。

\section{刷写、设备与可选 CBO}

\texttt{pmdev\_write\_page} 后调用 \texttt{pmem\_persist}。默认 flush 逐 64\,B 行 touch 并计数；\texttt{pmem\_fence} 含编译器屏障与内存 fence。可选 \texttt{PMEM\_ZICBOM} 发出 \texttt{CBO.CLEAN}（假定块大小与权限）。实验性 NVDIMM：补丁 QEMU 注册独立设备，内核在 \texttt{kalloc} 前解析私有 FDT 并校验地址/大小。

\section{维护：页复用与显式压缩}

可达性扫描覆盖两槽根与 run 链；不可达页可分配给新元数据或新 run。\texttt{pmcompact} 要求 memtable/stream 空，预留不相交暂存后写出新链并一次发布；忙或空间不足返回明确错误。这不是在线全设备 GC。

\section{验收脚手架}

\begin{description}[leftmargin=1.5em,style=nextline]
  \item[旧自测] \texttt{ONBOX\_PMFS\_*} 阶段码；\texttt{contract-full} 等破坏性，仅 disposable。
  \item[重启 oracle] \texttt{make test-recovery*}：临时后端、guest 步骤、\texttt{SIGKILL}、同后端读回；覆盖挂载拒绝、文件/KV、目录/链接、插桩切点。
  \item[模型] \texttt{make test-model}：有界协议枚举 + 负控反例。
  \item[轨迹] 测试构建记录写范围/flush/fence，采样重放提前逐出等。
\end{description}

尖刺价值：正确性依赖「可脚本判定 + 诚实假设」，而非口头演示。

\chapter{实验与分析}
\label{chap:eval}

\section{实验环境}

\begin{table}[htbp]
  \centering
  \caption{实验环境}
  \label{tab:env}
  \begin{tabular}{ll}
    \toprule
    项 & 配置 \\
    \midrule
    基线 & xv6-riscv（课题分支） \\
    模拟器 & QEMU \texttt{virt}，\texttt{CPUS=1} \\
    后端 & EMUL（8\,MiB）/ MAP / PMEM2（见第~\ref{chap:design}~章） \\
    回归 & \texttt{usertests}、\texttt{contract}/\texttt{contract-strict}/MAP/PMEM2（CI） \\
    计时 & 内核 \texttt{r\_time()} 相对单位；每 $N$ 取 5 trial 中位数；每 $N$ 独立开机 \\
    \bottomrule
  \end{tabular}
\end{table}

数据复现：\texttt{make bench-sweep} 回填 \texttt{docs/BENCH.md}；契约：\texttt{make contract}、\texttt{contract-strict}、\texttt{contract-pmem}、\texttt{contract-pmem2}。公平消融见 \texttt{docs/ABLATION.md}。语义与故障矩阵见第~\ref{sec:semantics}~节与 \texttt{docs/SEMANTICS.md}。

\textbf{两套结论：}架构面——相对 xv6 日志路径的量级差；公平面——同载荷消融效应量。禁止把 \texttt{fratio} 读成同等 POSIX API 加速比。

\section{功能、崩溃与掉电（尖刺证据）}

\begin{table}[htbp]
  \centering
  \caption{契约验收用例}
  \label{tab:contract}
  \begin{tabular}{llll}
    \toprule
    用例 & 命令/阶段 & 期望 & 说明 \\
    \midrule
    总闸 & \texttt{make contract} & \texttt{ONBOX\_PMFS\_PASS} & EMUL 必过子集 \\
    严格 & \texttt{contract-strict} & 同上 + A0--A4/\texttt{FAULT\_*} & POS/NEG/效应量 \\
    MAP & \texttt{contract-pmem} & \texttt{mode=MAP}+PMEM1 & 窗口映射 \\
    掉电 & \texttt{contract-pmem2} & \texttt{PMEM2\_PASS} & 文件后端杀启 \\
    remount & \texttt{REMOUNT} & \texttt{equal=1} & sync 后字节等价 \\
    故障 & \texttt{FAULT\_*} & 不可见 / chk 拒绝 & 未 sync；半页撕坏 \\
    POSIX & \texttt{pmfs\_posix} & \texttt{POSIX\_PASS} & 真实文件 API，不升主线 \\
    \bottomrule
  \end{tabular}
\end{table}

分析：契约套件使「完好 SB / 未提交不可见 / remount 等价 / 半页拒绝 / 窗口 persist」可机器判定；相对仅演示双槽结构，这是方法学差异，构成本文尖刺。

\section{性能：随 $N$ 变化的延迟}

对比路径故意不对等（表~\ref{tab:unequal}）。ratio$\gg 1$ 主要反映避免 virtio 日志与块栈的架构差，\textbf{不可}解读为同等完整 POSIX 下的千倍加速。

\begin{table}[htbp]
  \centering
  \caption{故意不对等的对比维度}
  \label{tab:unequal}
  \begin{tabular}{p{2.0cm}p{5.0cm}p{5.2cm}}
    \toprule
    维度 & PMFS-KV / File & xv6 fs+log \\
    \midrule
    API & KV：put$\times N$+seal+get；
         File：create/write/fsync$\times N$
         & open/create+write(8B)$\times N$+sync \\
    持久语义 & 双槽 SB + inode/dentry 页 & 块日志检查点 \\
    I/O 栈 & 字节寻址窗口 / EMUL & virtio-blk + bio + log \\
    元数据 & 有限 inode/dentry（无块位图日志） & 目录+inode+位图+log \\
    \bottomrule
  \end{tabular}
\end{table}

\begin{table}[htbp]
  \centering
  \caption{EMUL 下随 $N$ 变化的中位数与相对吞吐（5 trial，CDC 融合后）}
  \label{tab:bench-n}
  \begin{tabular}{rrrrrrr}
    \toprule
    $N$ & med\_kv & med\_file & med\_xv6 & fratio$\times$100 & thr\_file & wa\_x100 \\
    \midrule
    8  & 33507 & 304451 & 5684939  & 1867 & 26 & 15384 \\
    16 & 38448 & 486211 & 11208338 & 2305 & 32 & 17889 \\
    32 & 53065 & 659056 & 25648605 & 3891 & 48 & 17721 \\
    \bottomrule
  \end{tabular}
\end{table}

\FigNLatency

观察：（1）xv6 随 $N$ 近似线性；（2）PMFS-File 每文件 fsync 仍避开块日志栈，fratio$\gg 1$；（3）File 路径因 CDC/stream/meta 提交高于纯 KV，但仍与 xv6 差一个数量级以上；（4）结论写作\textbf{量级优势 + 边界}。

\section{写放大}

\FigWaN

小 value（8\,B）相对 4\,KiB 页对齐主导放大。CDC 融合后本负载下 $\mathrm{wa\_x100}$ 维持在高位且随 $N$ 近似平台（表~\ref{tab:bench-n}），属页粒度与 run 打包的预期行为；与 xv6 日志放大不同源，不可横向比绝对值。

\section{MAP 与 EMUL 对照}

\begin{table}[htbp]
  \centering
  \caption{$N{=}32$ 时 EMUL 与 MAP}
  \label{tab:map-emul}
  \begin{tabular}{lrrrr}
    \toprule
    mode & med\_kv & med\_file & med\_xv6 & fratio$\times$100 \\
    \midrule
    EMUL & 53065  & 659056  & 25648605 & 3891 \\
    MAP  & 113804 & 4301198 & 24439543 & 568 \\
    \bottomrule
  \end{tabular}
\end{table}

\FigMapEmul

\FigThr

文件路径相对吞吐在 EMUL 下稳定高于 xv6 一个数量级以上（表中 \texttt{thr\_file}）。MAP 抬升软件路径开销后 fratio 下降但仍 $\gg 1$。EMUL 仍\textbf{不}代表真 NVDIMM 延迟。

\section{公平消融（公平面）}

\begin{table}[htbp]
  \centering
  \caption{A0--A4 公平消融效应量（EMUL，\texttt{contract-strict} 实测）}
  \label{tab:ablation}
  \begin{tabular}{lccrrr}
    \toprule
    变体 & CDC & batch seal & cycles & wa\_x100 & 解读 \\
    \midrule
    A0 & 开 & 开 & 30029 & 1600 & 默认策略 \\
    A1 & 关 & 开 & 27490 & 1600 & 定长+批量 \\
    A2 & 开 & 关 & 52066 & 4800 & 每 chunk seal，WA/时延升 \\
    A3 & 关 & 关 & 36238 & 3200 & 定长+立即 seal \\
    A4 & 开 & 开 & 24146 & 1600 & A0 复测（稳定性） \\
    \bottomrule
  \end{tabular}
\end{table}

同 payload 下，关闭批量 seal（A2/A3）显著抬升 $\mathrm{wa\_x100}$ 与 cycles；A0/A4 同配置复测同量级，说明效应量可重复。公平面只据此表陈述，\textbf{不}与表~\ref{tab:bench-n} 的 fratio 混读。

\section{实现难点与失败史}

下列踩坑用于证明系统是「调出来的」，并可复核。

\paragraph{P0 启动栈溢出}
\texttt{main()} 跑在原 4\,KiB/hart 的 \texttt{stack0} 上；自测栈上放置多页缓冲导致溢出，污染邻近内存。修复：页缓冲改 \texttt{static}；\texttt{stack0} 扩至 16\,KiB/hart。

\paragraph{P1 \texttt{rdcycle} 非法指令}
\texttt{timerinit} 未开放 \texttt{CY} 时，S 态 \texttt{rdcycle} 触发 \texttt{kerneltrap}。修复：改用已允许的 \texttt{r\_time()}，并在论文中不作墙钟 $\mu$s 解读。

\paragraph{P2 QEMU NUMA / RAM 尺寸不匹配}
仅配置 160\,M 文件后端而默认 $-m$ 仍为 128\,M 时，QEMU 立即退出，契约表现为超时。修复：文件后端启动同步 $-m$\,160M。

\paragraph{P3 MAP 在 $-m$\,128M 下写 PMEMBASE}
无物理后备窗口上的写探测导致 store fault。修复：仅 \texttt{-DPMEM\_MAP} 时映射并探测。

\paragraph{P4 Bench 过早杀进程截断契约行}
匹配到 \texttt{BENCH\_PASS} 子串即杀 QEMU，导致字段截断。修复：等待 \texttt{median\_xv6=} / \texttt{phys=} / \texttt{wa\_x100=} 齐全。

\paragraph{P5 多 trial 撑爆段数}
连续 seal 耗尽段视图导致后续 trial 失败。修复：提高 run/段余量、attach 压力阀、\textbf{每 $N$ 独立开机}扫参。

\paragraph{P6 题面「文件系统」与实现漂移}
早期曾偏并行 KV。修复：升格 dir/extent/CDC + 深化 \texttt{pmfs\_posix}，并以语义清单写死差集，避免题实不符。

\section{局限与威胁效度}

\begin{enumerate}[leftmargin=2em]
  \item QEMU 非真 NVDIMM；附录给出扩展路径但不喧宾夺主。
  \item 架构面性能对比 API 异构；结论限于量级说明。公平面仅 A0--A4 同载荷可比。
  \item 文件语义为可测子集：无 \texttt{lseek}/\texttt{mmap}/\texttt{rename}；容量受 extent/chunk 上限约束；默认 \texttt{CPUS=1}。
  \item 计时为相对单位；外推墙钟或真机延迟需重新校准。
\end{enumerate}

\section{本章小结}

尖刺（契约套件，含 remount/故障）+ 支撑柱（文件子系统/KV 核/PMEM2）+ 两套结论口径 + 可复核踩坑，构成实验章叙事主线；性能图服务于量级与写放大机制说明，消融表服务于公平面效应量，而非绝对公平微基准夺冠。

\chapter{结论与展望}
\label{chap:concl}

\section{工作总结}

本文在 xv6-riscv 上设计并实现了并行于根文件系统的 \textbf{PMFS v3}：固定挂载点 \texttt{/pm} 上的受限文件语义、COW 元数据与不可变 KV run、双槽 checkpoint 提交、行级 flush 请求与可选 Zicbom，以及面向 QEMU 的实验性独立 NVDIMM 补丁路径。系统以进程重启切点 oracle、有界规格负控与轨迹采样验收提交行为，并以对齐工作量基准报告相对原生 \texttt{fs+log} 的软件路径延迟/吞吐量级差。

主要结论：

\begin{enumerate}[leftmargin=2em]
  \item \textbf{方法学：}不可替代点是可复现的验收资产与诚实假设清单，而非又一个双槽结构或完整生产 FS。
  \item \textbf{系统：}成功 \texttt{fsync} 后可读回完整版本；空白/损坏拒绝挂载；PM 路径不经块日志栈。
  \item \textbf{实验：}架构面展示量级优势但持久契约不等价；公平面展示软件请求效应，二者不可混读。
\end{enumerate}

对照原题四要件：NVDIMM 映射在\textbf{实验补丁路径}有条件达标；去块布局实质达标；刷写一致性在\textbf{声明假设}下达标而物理掉电未证；架构级性能叙事在软件路径上成立。

\section{不足}

真机持久域校准、上游通用设备发现、完整 POSIX、可扩展并发、近满可持续 GC、以及论文正式模板化定稿仍不足。旧 PDF/历史不等量基准不得作为新结论。

\section{展望}

\begin{itemize}[leftmargin=2em]
  \item 与导师书面确认设备模型与并行 \texttt{/pm} 产品形态；
  \item 将实际 store/flush 轨迹更紧地接入状态空间检查；
  \item 在保持尖刺的前提下扩展可测语义边界与论文图表重绘。
\end{itemize}

答辩一句话：\textbf{在教学 OS 上交付可经普通文件 API 访问的 PM 文件子系统，并以可脚本化故障验收与诚实假设证明提交协议；相对原生块日志展示架构级软件路径优势，但不声称无条件物理掉电或真机持久加速。}


\bibliographystyle{plain}
\bibliography{refs}

\appendix
\chapter{复现命令与设备说明}
\label{app:repro}

\section{安全警告}

\texttt{contract*}、\texttt{test-recovery*}、\texttt{bench-sweep}、部分 \texttt{test-xv6.py} 目标会 \texttt{make clean} 并重建 \texttt{fs.img}；boot 自测与 \texttt{pmformat --force} 会破坏介质。请仅在可丢弃 checkout 运行并先备份。

\section{推荐验收命令}

\begin{lstlisting}
make test-model
CPUS=3 make test-vfs
make test-recovery CPUS=3 RECOVERY_ARGS="--log-dir .ci-logs/recovery"
make test-recovery-zicbom CPUS=3 RECOVERY_ARGS="--log-dir .ci-logs/zicbom"
# 补丁 QEMU 后：
# python3 tools/check_nvdimm_device.py --qemu /path/to/qemu-system-riscv64
# make test-recovery-nvdimm QEMU=/path/to/qemu-system-riscv64
make contract-full          # 破坏性旧自测
python3 tools/bench_perf.py --runs 10 --warmup 1 --samples 32 \
  --output /tmp/pmfs-matched
\end{lstlisting}

文档索引见仓库 \texttt{docs/README.md}（含语义、契约、QEMU/NVDIMM 入口）。

\section{关于「NVDIMM」与「掉电」}

\begin{itemize}[leftmargin=2em]
  \item 未打补丁的 MAP / 整机共享 RAM \textbf{不是} NVDIMM 设备。
  \item 实验补丁提供独立 \texttt{nvdimm} 冷插拔与私有 FDT，可经 QMP 核验；非上游通用、无 NFIT/热插拔。
  \item 默认 \texttt{pmem\_flush} 非硬件写回；可选 CBO 在 TCG 下仍无真实持久域。
  \item 进程 \texttt{SIGKILL} 验证的是该故障模型下的读回，\textbf{不是}物理掉电。
\end{itemize}

\section{论文 PDF 编译}

\begin{lstlisting}
cd thesis/latex
make          # xelatex + bibtex + xelatex x2
\end{lstlisting}

封面姓名/学院等宏在 \texttt{main.tex}；请替换为学校正式模板后再印刷提交。


\backmatter
\chapter*{致\quad 谢}
\addcontentsline{toc}{chapter}{致谢}
\markboth{致谢}{致谢}

感谢指导教师在选题、系统边界与论文表述上的悉心指导。感谢教学操作系统 xv6 与 RISC-V / QEMU 开源社区提供的可读代码与实验平台。感谢课程与实验室同学在讨论、复现与答辩准备中的帮助。

本课题将工程验收（可机器判定契约）置于与数据结构同等重要的位置；文中踩坑记录旨在便于后来者复核。文责自负，疏漏之处恳请批评指正。


\end{document}
